\chapter{Вектороподобный mass/edge-селектор изотипического канала}
\label{ch:v8-baryon-c0-isotypic-channel-vectorlike-mass-edge-selector}

\section{Два поля на общем правом триплете}

После аномально безопасного завершения имеются два прямоугольных поля с
общим правым пространством:
\begin{equation}
 B\in\operatorname{Hom}(W_{eXZ},V_{LY})\simeq M_{2\times3}(\mathbb C),
 \qquad
 M\in\operatorname{Hom}(W_{eXZ},V_{XZ})\simeq M_{2\times3}(\mathbb C).
 \label{eq:v8-isotypic-mass-selector-two-fields}
\end{equation}
Первое поле должно иметь ранг один и норму три; второе --- ранг два и
одномерное правое ядро. Если оба требования задаются независимо, их правые
направления не обязаны совпасть.

\section{Минимальный выровненный подродитель}

Рассмотрим неотрицательный функционал
\begin{equation}
 \begin{split}
 \mathcal S_{BM}(B,M)={}&
 \bigl(\Tr BB^\dagger-3\bigr)^2
 +4\det(BB^\dagger)\\
 &+\|MM^\dagger-I_2\|_F^2
 +\|MB^\dagger\|_F^2.
 \end{split}
 \label{eq:v8-isotypic-mass-selector-action}
\end{equation}
Первые два слагаемых являются прежним coherence-потенциалом. Третье делает
$M$ частичной изометрией ранга два, а последнее связывает её ядро с
правой линией $B$.

Поскольку все четыре слагаемых неотрицательны,
\begin{equation}
 \mathcal S_{BM}=0
 \Longleftrightarrow
 \begin{cases}
 \Tr BB^\dagger=3,\quad \rank B=1,\\
 MM^\dagger=I_2,\quad \rank M=2,\\
 MB^\dagger=0.
 \end{cases}
 \label{eq:v8-isotypic-mass-selector-zero-set}
\end{equation}
Если
\begin{equation}
 P_B=\frac13B^\dagger B,
 \qquad
 P_M=M^\dagger M,
 \label{eq:v8-isotypic-mass-selector-projectors}
\end{equation}
то на нулевом множестве
\begin{equation}
 P_B^2=P_B,quad \rank P_B=1,
 \qquad
 P_M^2=P_M,quad \rank P_M=2,
 \qquad
 P_B+P_M=I_3.
 \label{eq:v8-isotypic-mass-selector-complementarity}
\end{equation}
Следовательно,
\begin{equation}
 \ker M=\operatorname{im}P_B,
 \label{eq:v8-isotypic-mass-selector-kernel-alignment}
\end{equation}
то есть coherence-конденсат сам выбирает лёгкое массовое направление.

\section{Ковариантность без фиксированного spurion}

Функционал инвариантен относительно
\begin{equation}
 B\mapsto g_BBh^\dagger,
 \qquad
 M\mapsto g_MMh^\dagger,
 \qquad
 (g_B,g_M,h)\in U(2)\times U(2)\times U(3).
 \label{eq:v8-isotypic-mass-selector-covariance}
\end{equation}
Поэтому ненулевая масса не вводится фиксированным $SO(3)$-ломающим
интертвейнером: направление появляется как совместный вакуум двух полей.
Такой norm-square язык согласуется с моментным подходом к представлениям
колчанов \cite{v8-isotypic-mass-selector-quiver} и с появлением новых
скалярных связей при вектороподобных расширениях спектральной геометрии
\cite{v8-isotypic-mass-selector-scalar}.

Явный нулевой представитель равен
\begin{equation}
 B_0=\begin{pmatrix}\sqrt3&0&0\\0&0&0\end{pmatrix},
 \qquad
 M_0=\begin{pmatrix}0&1&0\\0&0&1\end{pmatrix}.
 \label{eq:v8-isotypic-mass-selector-vacuum}
\end{equation}
На вещественном двенадцатимерном срезе его полный гессиан имеет
\begin{equation}
 \Spec\Hess_{(B_0,M_0)}\mathcal S_{BM}
 =\{0\}^{\times4}\cup\{8\}^{\times5}
 \cup\{24\}\cup\{26\}^{\times2}.
 \label{eq:v8-isotypic-mass-selector-hessian}
\end{equation}
Четыре нуля касательны к орбите; отрицательных поперечных мод нет.

\section{Почему абсолютный канал не выбран}

Для любой общей правой замены базиса $h$ пара
\begin{equation}
 (B_0,M_0)\longmapsto(B_0h^\dagger,M_0h^\dagger)
 \label{eq:v8-isotypic-mass-selector-vacuum-orbit}
\end{equation}
остаётся нулевой. В частности, циклическая перестановка трёх каналов даёт
другой проектор $P_B$ с той же энергией. Поэтому parent выравнивает ядро и
конденсат, но не выбирает, какая комбинация является физической family-line.

\section{Полный граф остаётся вне действия}

Функционал использует выбранные прямоугольники $B$ и $M$, тогда как строгий
граф предыдущего гейта разрешает шесть дополнительных комплексных полей
\begin{equation}
 z_{\rm extra}in\mathbb C^6.
 \label{eq:v8-isotypic-mass-selector-extra-fields}
\end{equation}
При буквальном продолжении
\begin{equation}
 \mathcal S_{\rm ext}(B,M,z_{\rm extra})=\mathcal S_{BM}(B,M)
 \label{eq:v8-isotypic-mass-selector-flat-extension}
\end{equation}
они дают двенадцать новых вещественных нулей. Сигнатура проверенного
вещественного среза полного продолжения равна
\begin{equation}
 (n_-,n_0,n_+)=(0,16,8).
 \label{eq:v8-isotypic-mass-selector-full-signature}
\end{equation}
Тем самым построен точный aligned-subparent, но projector на выбранные
рёбра уже предположен в его области определения.

Реестры имеют вид
\begin{equation}
 \begin{aligned}
 N_{\rm aligned\ shape}&=8/8,
 &N_{\rm full\ graph\ selector}&=0/5,\\
 N_{\rm parent\ origin}&=0/4,
 &\mathcal M_{\rm new}&=\{M,c_i,\Pi_E,SO(3)_{\rm diag}\}.
 \end{aligned}
 \label{eq:v8-isotypic-mass-selector-ledgers}
\end{equation}

\begin{theorem}[Точное выравнивание массового ядра]
Неотрицательный функционал $\mathcal S_{BM}$ имеет нулевое множество, на
котором coherence-проектор ранга один и массовый проектор ранга два взаимно
ортогональны и в сумме дают единицу. Поэтому лёгкое ядро массы совпадает с
правой линией coherence-конденсата без фиксированного направляющего spurion.
В проверенном вещественном срезе вакуум поперечно устойчив.
\end{theorem}

\begin{proposition}[Отсутствие coordinate edge-селектора]
Совместный parent сохраняет общую правую $U(3)$-орбиту и не выбирает
абсолютное канальное направление. Кроме того, шесть первопорядково
разрешённых полей не входят в его действие и остаются двенадцатью
вещественными плоскими модами. Следовательно, полный строгий граф этим
функционалом не выбран.
\end{proposition}

\begin{nogocriterion}[Запрет скрытого edge-проектора]
Нельзя считать условие $MB^\dagger=0$ полным физическим селектором, если
поля $B$ и $M$ заранее определены только на желаемых блоках. Требуется
вложить тот же механизм во все девять новых рёбер и получить положительную
кривизну на шести конкурирующих направлениях без ручного зануления.
\end{nogocriterion}

\begin{protocol}\begin{enumerate}
\item общий правый изотипический триплет: \textbf{построен};
\item coherence-поле ранга один: \textbf{получено};
\item массовое поле ранга два: \textbf{получено};
\item точное условие выравнивания: \textbf{$MB^\dagger=0$};
\item проекторное тождество: \textbf{$P_B+P_M=I_3$};
\item фиксированный $SO(3)$-ломающий spurion: \textbf{не требуется};
\item отрицательные моды вещественного гессиана: \textbf{ноль};
\item aligned-shape ledger: \textbf{$8/8$};
\item абсолютное channel-направление: \textbf{не выбрано};
\item пропущенных разрешённых комплексных рёбер: \textbf{шесть};
\item дополнительных вещественных нулей: \textbf{двенадцать};
\item full-graph selector ledger: \textbf{$0/5$};
\item parent-origin ledger: \textbf{$0/4$};
\item portal parent-origin: \textbf{не закрыт};
\item следующий гейт: \textbf{полнографовое вложение выровненного parent}.
\end{enumerate}\end{protocol}

Машинный сертификат сохранён в
\path{s2t_v8_baryon_c0_singlet_triplet_central_gap_isotypic_channel_vectorlike_mass_edge_selector_gate_results.json}.

\begin{thebibliography}{9}
\bibitem{v8-isotypic-mass-selector-quiver}
M.~Harada, G.~Wilkin,
\emph{Morse Theory of the Moment Map for Representations of Quivers},
arXiv:0807.4734.
\bibitem{v8-isotypic-mass-selector-scalar}
C.~A.~Stephan,
\emph{New Scalar Fields in Noncommutative Geometry},
arXiv:0901.4676.
\end{thebibliography}